% ============================================================
% intro.tex
% ============================================================

This chapter motivates the research question---what determines the
timing and type of resolution in federal securities class
actions---situates the thesis within the empirical litigation
literature, and summarizes the principal contributions.

Federal securities class actions represent one of the most
economically significant categories of civil litigation in the
United States, with annual filings fluctuating between 200 and
400 cases and aggregate settlements frequently exceeding \$2
billion \citep{cornerstone2024, nera2024}. 

These cases arise when investors allege violations of federal
securities laws---primarily Section~10(b) of the Securities
Exchange Act of 1934 and SEC Rule 10b-5, which prohibit
fraudulent misrepresentations, or Section~11 of the Securities
Act of 1933, which imposes strict liability for material
misstatements in registration statements. Defendants include
publicly traded corporations and their officers, with plaintiffs
ranging from individual retail investors to large institutional
investors who increasingly serve as lead plaintiffs following the
Private Securities Litigation Reform Act of 1995 (PSLRA)
\citep{cox2008, cheng2010}.

The stakes are substantial. For defendants, direct legal
expenses and settlement payouts range from millions to over a
billion dollars in so-called ``mega-settlements''
\citep{cornerstone2024}. Beyond monetary costs, securities
litigation generates elevated cost of capital, shifts in
investment policies, managerial distraction, and reputational
damage \citep{karpoff2008, graham2008}. For plaintiffs and their
attorneys, duration and outcome directly affect net recovery and
the feasibility of pursuing future cases. For courts and
policymakers, resolution patterns inform debates over the
efficacy of procedural reforms and the balance between deterring
fraud and curbing frivolous litigation.

Despite the prominence of securities class actions, empirical
research on litigation duration and outcomes has been limited in
both methodology and data quality. Practitioner reports provide
valuable descriptive statistics, but do not employ formal survival
analysis or competing-risks frameworks \citep{cornerstone2024,
nera2024}. Academic studies have applied Cox proportional hazards 
models to settlement speed in selected samples \citep{brochet2014}, 
but ignore the competing nature of case dismissal. In practice, 
nearly all core federal securities class actions ultimately either 
settle or are dismissed. Treating these mutually exclusive outcomes
independently introduces selection bias and provides an incomplete 
picture of litigation dynamics. Furthermore, this methodological gap 
is compounded by a profound data challenge: administrative court 
databases frequently misclassify voluntary dismissals---cases
in which the parties privately negotiate a settlement but the
court records only a dismissal---masking these ``hidden
settlements'' and distorting true resolution probabilities
\citep{scales2024}.

This thesis addresses these gaps along two dimensions. Methodologically, it applies, to our knowledge, the first full competing-risks survival framework to federal securities class actions at population scale, drawing on 12,968 cases from the Federal Judicial Center's (FJC) Integrated Database (IDB) spanning 1990--2024. We characterize litigation duration with Aalen-Johansen cumulative incidence functions (CIFs), estimate both cause-specific Cox proportional hazards and Fine-Gray subdistribution models, which directly parameterize the cumulative probability of each outcome, and use inverse probability of treatment weighting (IPTW) as a composition-adjustment strategy rather than a causal design. Circuit-level shared frailty models complement these estimates by quantifying residual jurisdictional heterogeneity.

Substantively, the results show that the PSLRA's association with case outcomes is not uniform over time. Its strongest and most defensible association is an increase in early dismissal hazard that fades over litigation duration. The settlement evidence points in the same direction in several specifications, including the piecewise model, but is materially more fragile under circuit-specific, time-trend, and cluster-robust checks. The thesis also shows that judicial geography is a major structural determinant of outcomes and that MDL consolidation changes the relationship between instantaneous hazards and cumulative outcome probabilities.

On the data side, the thesis resolves two measurement
problems in the FJC Integrated Database \citep{clermont2002,
alexander1991}: ambiguity between settlement-like dismissals
and true dismissals, and ambiguity within judgment-bearing
dispositions. Chapter~\ref{ch:data} develops those coding
choices and their sensitivity analyses in detail. Because the
placebo and time-trend checks leave temporal confounding
unresolved, the thesis frames its estimates as associational
or composition-adjusted rather than causally identified.

The remainder of this thesis is organized as follows.
Chapter~\ref{ch:litreview} reviews the empirical securities
litigation literature and the survival analysis methods that
motivate our approach. Chapter~\ref{ch:methodology} develops
the competing-risks framework, propensity-score weighting
procedure, and shared frailty specification.
Chapter~\ref{ch:data} describes the IDB data, sample
construction, and disposition coding schemes.
Chapter~\ref{ch:results} presents the results.
Chapter~\ref{ch:discussion} interprets the findings,
discusses limitations, and situates the contributions within
the broader literature. Chapter~\ref{ch:conclusion} concludes
the thesis and outlines directions for future research.
